\originalsection{18.04期末模拟考试}
\sourceinfo{Jeremy Orloff, Spring 2018. 题面 practice-final.pdf, PDF p.1; 官方解答 practice-final-qa, pp.1--2. MIT OCW, CC BY-NC-SA 4.0. 3道大题. 原卷说明实际期末考试提供所附 Laplace 变换表; 本公开模拟题未考该表, 不能据此推断完整期末范围.}
\begin{exercise}\label{ass:pfinal-1}
\textbf{原题 PFinal.1。} 判断是否在全复平面亚纯: (a) $z^5$; (b) $z^{5/2}$; (c) $\ee^{1/z}$; (d) $1/\sin z$.
\end{exercise}
\begin{solution}据官方解答翻译. (a) 是, 多项式整. (b) 否, 绕0延拓会改变符号, 不能在全平面定义单值亚纯函数. (c) 否, 0是本性奇点. (d) 是, $\sin z$ 仅在 $k\pi$ ($k\in\Z$) 有简单零, 倒数在这些点有简单极点, 无其他有限奇点. “全平面亚纯”不要求无穷远也亚纯.\end{solution}
\begin{exercise}\label{ass:pfinal-2}
\textbf{原题 PFinal.2。} (a) 对 $f(z)=(z-2)^2z^3/((z+5)^3(z+1)^3(z-1)^4)$, 求正向 $\int_{|z|=3}f'/f\dd z$. (b) 求 $6z^4+z^3-2z^2+z-1$ 在单位圆盘内零点数. (c) $f$ 在单位闭圆盘的邻域上解析且边界 $|f|<1$, 证明内部恰有一个不动点. (d) 判断: 在简单闭曲线及内部解析的函数 $f$, 若内部有 $n$ 个零点, 则 $f'$ 内部有 $n-1$ 个零点.
\end{exercise}
\begin{solution}
据官方解答翻译并明确按重数计数. (a) 内部零点总阶数 $2+3=5$, 极点总阶数 $3+4=7$, $-5$ 在外. 辐角原理给 $2\pi\ii(5-7)=-4\pi\ii$. (b) 单位圆上余项模不超过5, 严格小于 $|6z^4|=6$. Rouché 定理给4个零点, 边界也无零点.

(c) 圆周上 $|f(z)|<|-z|=1$, 故 $f(z)-z$ 与 $-z$ 内部零点数相同, 均为1. 这还说明不动点是单根. 本题与18.112历史期末第6题同题, 两处保留编号与出处. (d) 假. 取 $f(z)=\ee^z-1$, 圆 $|z|=3\pi$ 内有 $0,\pm2\pi\ii$ 三个简单零点, 边界无零; $f'=\ee^z$ 无零点. 多项式在某些条件下的结果不能无条件推广到解析函数与任意曲线.
\end{solution}
\begin{exercise}\label{ass:pfinal-3}
\textbf{原题 PFinal.3。} 令 $A=\{z:0\le\Re z\le\pi/2,\Im z\ge0\}$, $B$ 为第一象限, 证明 $f(z)=\sin z$ 把 $A$ 保角映到 $B$.
\end{exercise}
\begin{solution}
据官方解答翻译边界计算; 完整单射、满射论证为编者补充（AI）. 原文 $A$ 闭而“保角”定义于开域, 应分为内部保角双射和边界连续对应: $A^\circ=\{0<x<\pi/2,y>0\}$ 映到开第一象限 $B^\circ$, 闭 $A$ 映到闭第一象限. 在角点 $\pi/2$ 导数为0, 不声称边界处保角.

边界三段的像分别为 $\sin(\ii y)=\ii\sinh y$（正虚轴）, $\sin x$ ($0\le x\le\pi/2$, 实区间 $[0,1]$), $\sin(\pi/2+\ii y)=\cosh y$（实射线 $[1,\infty)$）. 内部的实部 $\sin x\cosh y>0$, 虚部 $\cos x\sinh y>0$.

单射由 $\sin z=\sin w$ 的全部可能 $z-w\in2\pi\Z$ 或 $z+w\in\pi+2\pi\Z$ 得到: 在内部实部差介于 $-\pi/2,\pi/2$, 前者仅 $z=w$; 后者虚部和正而右侧实, 不可能. 导数 $\cos z$ 在内部不为0.

满射可直接构造. 给 $a,b>0$, 寻找 $t=\sinh^2y>0$, 解 $a^2/(1+t)+b^2/t=1$. 左边在 $(0,\infty)$ 连续严格递减, 从无穷大降至0, 故唯一 $t$. 取 $y=\operatorname{arsinh}\sqrt t>0$, 定义 $\sin x=a/\sqrt{1+t},\cos x=b/\sqrt t$; 二者正且平方和1, 给唯一 $0<x<\pi/2$. 于是 $\sin(x+\ii y)=a+\ii b$. 只画边界并检验一个内点不足以替代这一论证.
\begin{center}\begin{tikzpicture}[scale=.85,>=Stealth]
\draw[->](-.3,0)--(2.4,0) node[right]{$x$};\draw[->](0,-.3)--(0,2.8) node[above]{$y$};\fill[main!10](0,0)rectangle(1.57,2.5);\draw[thick](0,0)--(0,2.5);\draw[thick](0,0)--(1.57,0)--(1.57,2.5);\node[below]at(1.57,0){$\pi/2$};\node at(.8,1.5){$A^\circ$};
\draw[->](2.7,1.3)--(4.1,1.3) node[midway,above]{$\sin$};
\begin{scope}[xshift=5cm]\fill[main!10](0,0)rectangle(2.4,2.5);\draw[->](-.3,0)--(2.9,0) node[right]{$\Re w$};\draw[->](0,-.3)--(0,2.8) node[above]{$\Im w$};\fill(1,0)circle(1.5pt) node[below]{$1$};\node at(1.2,1.5){$B^\circ$};\end{scope}
\end{tikzpicture}\end{center}
\end{solution}
